%%-*-latex-*-

\documentclass[11pt, a4paper]{article}

\usepackage[utf8]{inputenc}
\usepackage[T1]{fontenc}       % Required for \DJ and hyphenation
\usepackage[french]{babel}
\usepackage{amsmath}
\usepackage{graphicx}

\title{Examen de programmation fonctionnelle en Objective Caml}
\author{Christian Rinderknecht}
\date{Mardi 22 avril 2003}

\begin{document}

\maketitle

Les exercices peuvent être traités de façon indépendante.
Barême: A:3, B:3, C:10, D:4.

\smallskip

\textbf{Exercice A:}
Donner le type de chacune des expressions suivantes:

\noindent
\textbf{A.1}: \verb"fun (x,y) -> (y+1, not x)"\\
\textbf{A.2}: \verb"fun p -> (snd p, fst p)"\\
\textbf{A.3}: \verb"fun f -> fun (x,y) -> f(x+1) + f y"\\
\textbf{A.4}: \verb"fun f -> fun (x,y) -> (f(x+1), f y)"\\
\textbf{A.5}: \verb"fun f -> fun (x,y) -> (f x, f y)"\\
\textbf{A.6}: \verb"fun f -> fun (x,y) -> (f x + 1, f y)"
%A.7: \verb"fun (f,g) -> fun x -> (f x, g x)"\\

\medskip

\textbf{Exercice B:}
Pour chacun des types suivants, trouver une expression Caml
possédant ce type:

\noindent
%B.1: \verb"(int * int) -> bool"\\
\textbf{B.1}: \verb"int -> bool -> int"\\
\textbf{B.2}: \verb"int -> (bool * int)"\\
\textbf{B.3}: \verb"int list -> bool list"\\
\textbf{B.4}: \verb"('a * 'b) list -> 'a list"\\
\textbf{B.5}: \verb"('a * 'b) list -> 'b -> 'a"\\
\textbf{B.6}: \verb"'a list -> 'b list -> ('a * 'b) list"
q
\medskip

\textbf{Exercice C:}

Un arbre peut représenter une collection de mots de sorte qu'il soit
très efficace de vérifier si une séquence donnée de
caractères est un mot valide ou pas. Dans ce type d'arbre, appelé
un {\em{trie}}, les arcs ont une lettre associée, chaque nœud
possède une indication si la lettre de l'arc entrant est une fin de
mot et la suite des mots partageant le même début. La figure
suivante montre le \emph{trie} des mots «~le~», «~les~», «~la~» et
«~son~», l'astérisque marquant la fin d'un mot:
\begin{center}
  \includegraphics[scale=0.8]{trie.eps}
\end{center}
\noindent
On utilisera le type Caml suivant pour implanter les {\em{tries}}:

\begin{verbatim}
type trie =
  { mot_complet : string option; suite : (char * trie) list }
\end{verbatim}

Le type prédéfini \verb+type 'a option = None | Some of 'a+ sert Ã
représenter une valeur éventuellement absente. Chaque nœud d'un
\emph{trie} contient les informations suivantes:

\begin{itemize}
  \item Si le nœud marque la fin d'un mot {\tt s} (ce qui
        correspond aux nœuds étoilés de la figure), alors le champ
        {\tt mot\_complet} contient {\tt Some s}, sinon ce champ
        contient {\tt None}.

  \item Le champ {\tt suite} contient une liste qui associe les
        caractères aux nœuds.
\end{itemize}

\medskip

\noindent
\textbf{C.1}: Écrire la valeur Caml de type \verb+trie+ correspondant
à la figure ci-dessus.\\
\textbf{C.2}: Écrire une fonction \texttt{compte\_mots} qui compte le
nombre de mots dans un \emph{trie}.\\
\textbf{C.3}: Écrire une fonction \texttt{select} qui prend en
argument un \emph{trie} et une lettre, et renvoie le \emph{trie}
correspondant aux mots commençant par cette lettre. Si ce \emph{trie}
n'existe pas (parce qu'aucun mot ne commence par cette lettre), la
fonction devra lancer une exception {\tt Absent}, que l'on
définira. On utilisera la fonction prédéfinie \texttt{List.assoc} pour
effectuer la recherche dans la liste des sous-arbres \verb+suite+.\\
\textbf{C.4}: Écrire une fonction \texttt{recherche} qui vérifie si
une cha\^{\i}ne de caractères est un mot dans un {\em trie}
donné. La fonction devra prendre un argument supplémentaire {\tt i},
qui représente la position dans le mot associée au nœud courant.
Le i\ieme{} caractère d'une chaîne {\tt s} s'obtient en écrivant {\tt
s.[i]}, le premier caractère étant numéroté $0$. La longueur de la
chaîne {\tt s} s'écrit {\tt String.length s}. On emploiera la fonction
\texttt{select}.

\medskip

\textbf{Exercice D:}

\noindent
On s'intéresse aux expressions booléennes écrites à l'aide des
connecteurs \texttt{or} («~ou~» booléen), \texttt{and} («~et~»
booléen), \texttt{not} (négation booléenne), des constantes
\texttt{true} et \texttt{false}, et de variables.  Par exemple: $$(x
~or~ y) ~and~ not(x ~and~ y)$$.

\noindent
\textbf{D.1}: Définir un type Caml \texttt{bool\_exp} pour ces
expressions.

\noindent
\textbf{D.2}: Écrire une fonction \texttt{eval} permettant d'évaluer
de telles expressions. Cette fonction devra prendre en paramètre un
environnement associant les variables à leur valeur.

\noindent
\textbf{D.3}:
Les connecteurs booléens considérés satisfont en particulier
les identités
\begin{align*}
not(a ~or~ b)  & = not(a) ~and~ not(b)\\
not(a~ and~ b) & = not(a)~ or~ not(b)\\
not(not(a))    & = a\\
not(true)      & = false\\
not(false)     & = true
\end{align*}

\noindent
Ces identités permettent de transformer toute expression booléenne
en une expression équivalente où les négations ne sont
appliquées qu'à des variables. Par exemple, l'expression $not(x
~and~ (y~ or~ not(z)))$ peut être transformée en $not(x) ~or~
not(y ~or ~not(z)))$ puis en $not(x) ~or~ (not(y)~ and ~not(not(z))))$
et enfin en $not(x) ~or~ (not(y) ~and~ z)$.

\noindent
Écrire une fonction Caml \texttt{normalise} qui réalise cette
transformation.

\end{document}
